\chapter{La fabbrica dei sogni e dei bias}

\begin{flushright}
	\textit{``Se vuoi controllare una popolazione, inizia dalle sue storie.''}\\[6pt]
	\textsc{Anonimo}
\end{flushright}

\bigskip
\bigskip

Se i capitoli precedenti hanno mostrato come i bias evolutivi si annidino nei meandri del nostro cervello antico, \`{e} ora di esplorare uno dei loro veicoli di diffusione pi\`{u} sofisticati: la cultura di massa. E non esiste laboratorio migliore per studiare questo fenomeno dell'impero Disney, la fabbrica di sogni che da quasi un secolo trasforma i nostri bias ancestrali in principesse cantanti e topolini parlanti. Non stiamo per rovinare la vostra infanzia. O forse s\`{i}, ma per una buona causa.

\section{Il topolino e il cavernicolo}

Disney non \`{e} soltanto un'azienda di intrattenimento: \`{e} un amplificatore culturale di bias cognitivi su scala planetaria.\footnote{Giroux, H. A. (1999). \textit{The Mouse that Roared: Disney and the End of Innocence}. Rowman e Littlefield, Lanham.} Ogni film, ogni personaggio iconico funziona come un cavallo di Troia che infiltra nelle menti infantili, e non solo, pattern di pensiero che risalgono alla savana africana. Il genio di Walt Disney \`{e} stato capire istintivamente qualcosa che gli psicologi evolutivi avrebbero formalizzato solo decenni dopo: il cervello \`{e} programmato per assorbire storie capaci di attivare circuiti neurali antichi.\footnote{Boyd, B. (2009). \textit{On the Origin of Stories: Evolution, Cognition, and Fiction}. Harvard University Press, Cambridge.} Non conta che quelle storie siano realistiche: conta che facciano vibrare le corde giuste del nostro hardware ereditato.

Prendiamo il concetto dell'eroe puro e senza macchia. Biancaneve, Cenerentola, Ariel sono tutte incarnazioni di quella che potremmo chiamare la sindrome del bravo bambino evolutivo: personaggi che non si arrabbiano mai, non cercano vendetta, accettano passivamente le ingiustizie e vincono attraverso la bont\`{a} e il sacrificio. \`{E} la trasfigurazione di antichissimi pattern di sottomissione sociale che servivano a mantenere coesione nel gruppo.\footnote{Haidt, J. (2012). \textit{The Righteous Mind: Why Good People Are Divided by Politics and Religion}. Vintage Books, New York.}

Ma per capire perch\'{e} questi personaggi ci colpiscano cos\`{i} in profondit\`{a}, indipendentemente dalla cultura in cui nasciamo, dobbiamo fare i conti con uno psicologo svizzero che, molto prima di Disney, aveva mappato il territorio su cui l'industria dell'intrattenimento avrebbe poi costruito il proprio impero. Carl Gustav Jung sosteneva che, sotto la superficie della nostra mente cosciente, esista un livello ancora pi\`{u} profondo del semplice inconscio personale: un \index[main]{inconscio collettivo}inconscio collettivo, condiviso da tutta la specie, popolato da figure ricorrenti che chiam\`{o} \index[main]{archetipi}archetipi.\footnote{Jung, C. G. (1954). \textit{Von den Wurzeln des Bewusstseins}. Rascher Verlag, Zurigo.} L'Eroe, la Grande Madre, il Vecchio Saggio, l'Ombra: figure che riaffiorano identiche nei miti di popoli che non hanno mai avuto contatto tra loro, perch\'{e} non sono invenzioni culturali ma predisposizioni della psiche umana, quasi organi mentali ereditati insieme al resto del nostro corredo biologico. Gli archetipi junghiani, letti con gli strumenti di questo libro, sono in fondo bias glorificati: scorciatoie cognitive antichissime, capaci di farci riconoscere all'istante un pattern narrativo, che la tradizione filosofica ha nobilitato con il nome di simbolo universale. E quando Disney riduce un personaggio a un solo tratto dominante, la principessa buona, la matrigna cattiva, l'eroe coraggioso, non sta inventando nulla di nuovo: sta semplicemente prelevando l'archetipo gi\`{a} pronto dal magazzino dell'inconscio collettivo e confezionandolo per la vendita al dettaglio, senza pi\`{u} lasciare spazio all'ambiguit\`{a} che, in Jung, ogni archetipo porta sempre con s\'{e}.

\section{La semplificazione morale: quando il mondo diventa cartone animato}

Disney non si limita a vendere personaggi: vende una visione del mondo radicalmente semplificata, che fa a pugni con qualunque comprensione sofisticata della natura umana. \`{E} quello che potremmo chiamare il bias manicheo, la tendenza a dividere il mondo in buoni assoluti e cattivi assoluti, senza sfumature intermedie. Non \`{e} una scelta casuale: \`{e} l'applicazione industriale di un bias cognitivo fondamentale, la nostra preferenza evolutiva per categorizzazioni nette e veloci.\footnote{Heath, C., e Heath, D. (2007). \textit{Made to Stick: Why Some Ideas Survive and Others Die}. Random House, New York.} Nella savana non c'era tempo per analisi psicologiche complesse del carattere del predatore che ci stava inseguendo: o era una minaccia da evitare, o non lo era. Disney ha trasformato questa necessit\`{a} evolutiva in un impero commerciale.

Il risultato sono generazioni cresciute con l'aspettativa inconscia che la realt\`{a} debba funzionare come un film Disney, con ruoli chiari, motivazioni trasparenti, finali dove il bene trionfa sempre. Quando poi si incontra la complessit\`{a} del mondo reale, dove i buoni a volte fanno cose terribili e i cattivi hanno motivazioni comprensibili, il cervello va in cortocircuito.

Questa operazione di levigatura non tocca solo i personaggi, ma la trama stessa delle storie da cui Disney attinge, e qui il discorso si fa pi\`{u} serio di quanto sembri a prima vista. Disney ha sistematicamente alterato le fiabe originali per renderle innocue e adatte a un pubblico di massa, eliminando elementi violenti, oscuri o inquietanti.\footnote{Zipes, J. (1995). \textit{Breaking the Magic Spell: Radical Theories of Folk and Fairy Tales}. University Press of Kentucky, Lexington.} La Cenerentola dei fratelli Grimm vedeva le sorellastre tagliarsi parti dei piedi per entrare nella scarpetta e infine essere accecate dai piccioni. La Biancaneve originale faceva torturare a morte la matrigna con scarpe di ferro roventi. Questa non \`{e} una questione di cattivo gusto vittoriano da correggere: \`{e} la distruzione di una funzione psicologica precisa, e a spiegarla meglio di chiunque altro fu lo psicoanalista Bruno Bettelheim ne \textit{Il mondo incantato}.\footnote{Bettelheim, B. (1976). \textit{The Uses of Enchantment: The Meaning and Importance of Fairy Tales}. Knopf, New York.} Per Bettelheim le fiabe tradizionali, con tutta la loro violenza apparentemente gratuita, funzionavano come una forma di terapia psicologica spontanea per l'infanzia: offrivano ai bambini, attraverso il filtro sicuro della finzione, un modo per elaborare paure reali e altrimenti indicibili, l'abbandono, la morte, la crudelt\`{a} degli adulti, la propria stessa rabbia distruttiva. Il lupo che divora la nonna, la matrigna punita, il gigante sconfitto non erano dettagli macabri aggiunti per gusto sadico: erano lo spazio simbolico in cui il terrore infantile, gi\`{a} presente comunque nella mente del bambino, trovava finalmente una forma, un nome, e quindi la possibilit\`{a} di essere superato. Sanitizzare la fiaba, per Bettelheim, non significa renderla pi\`{u} sicura: significa privare il bambino dello strumento che gli serviva per affrontare ci\`{o} che, dentro di lui, era gi\`{a} spaventoso comunque. Disney non ha reso le fiabe pi\`{u} innocenti: ha semplicemente tolto ai bambini la mappa per orientarsi nei propri stessi incubi.

\section{La repressione della rabbia}

Forse l'aspetto pi\`{u} insidioso della mitologia Disney \`{e} il trattamento riservato alla rabbia. Negli universi Disney questa emozione \`{e} quasi esclusivamente appannaggio dei villain: gli eroi positivi non si arrabbiano mai, nemmeno di fronte alle ingiustizie pi\`{u} flagranti. Nei miti antichi, da Achille a Mos\`{e}, la rabbia era invece riconosciuta come una forza potenzialmente distruttiva ma anche necessaria, l'emozione che permetteva di difendere i propri confini e di dire no quando serviva.\footnote{Campbell, J. (1949). \textit{The Hero with a Thousand Faces}. Pantheon Books, New York.} Disney ha sistematicamente estirpato questa dimensione dalle proprie narrazioni, sostituendola con una bont\`{a} passiva e remissiva. Il messaggio implicito \`{e} chiaro: i bravi bambini non si arrabbiano, sopportano, aspettano che qualcun altro, preferibilmente magico, risolva i loro problemi. \`{E} la perfetta preparazione psicologica per una societ\`{a} di consumatori docili.

Questa rimozione sistematica delle emozioni difficili, per\`{o}, non riguarda solo la rabbia: riguarda il modo stesso in cui una storia dovrebbe trattare qualunque emozione dolorosa, e su questo la filosofia occidentale aveva gi\`{a} una risposta pronta molti secoli prima di Walt Disney. Nella \textit{Poetica}, Aristotele spiegava che la tragedia teatrale ha un compito preciso: suscitare piet\`{a} e terrore nello spettatore per poi condurlo, attraverso la rappresentazione stessa di quelle emozioni, a una loro \index[main]{catarsi}catarsi, una purificazione.\footnote{Aristotele. \textit{Poetica}. Numerose edizioni.} Non si tratta di evitare il dolore, ma di attraversarlo in un contesto protetto, quello della finzione, cos\`{i} che l'emozione, vissuta e riconosciuta, possa infine sciogliersi invece di restare intrappolata. \`{E} esattamente il meccanismo che Bettelheim attribuiva alle fiabe originali, e che Disney interrompe sul nascere: se la rabbia non viene mai mostrata, se il dolore appare sempre gi\`{a} risolto o aggirato da un intervento magico, la catarsi aristotelica non ha modo di avvenire. Non si tratta pi\`{u} di purificare un'emozione attraversandola, ma di negarla in partenza, rimuoverla dallo schermo cos\`{i} come si vorrebbe rimuoverla dalla vita reale. Il pubblico Disney non piange e non si arrabbia insieme al personaggio per poi sentirsi alleggerito: resta semplicemente convinto, un film dopo l'altro, che certe emozioni non abbiano diritto di cittadinanza nemmeno nella finzione.

\section{Il controllo narrativo}

Se c'\`{e} un aspetto che distingue Disney da narratori pi\`{u} sofisticati, \`{e} la mania di controllo assoluto sull'esperienza dello spettatore. Dove altri autori lasciano spazio all'interpretazione, al simbolismo, al mistero, Disney imbocca letteralmente il pubblico, assicurandosi che ogni messaggio sia cristallino, ogni metafora esplicita, ogni significato gi\`{a} predigerito. Questo riflette e rinforza uno dei bias pi\`{u} pericolosi della modernit\`{a}: l'illusione di controllo. Disney vende il sogno che la vita possa essere prevista e addomesticata come una delle sue principesse, mentre la vita reale \`{e} piena di ambiguit\`{a} e significati che emergono lentamente, se emergono. L'effetto a lungo termine sono adulti che si rifugiano nelle certezze dell'infanzia Disney ogni volta che la realt\`{a} diventa troppo complessa da gestire, ed \`{e} forse anche per questo che l'impero ha costruito gran parte della propria fortuna recente su remake e sequel: vendere nostalgia per un'epoca in cui tutto sembrava sotto controllo, quando in realt\`{a} non lo era mai stato.

Uno degli aspetti pi\`{u} sottili di questa strategia \`{e} ci\`{o} che potremmo chiamare identificazione diretta: la creazione di personaggi deliberatamente generici, progettati come tele bianche su cui il pubblico possa proiettare se stesso, io voglio essere come la Bella, io voglio vivere come Elsa. Contrasta radicalmente con la funzione dei miti tradizionali, dove i personaggi non erano modelli da imitare ma archetipi, per tornare a Jung, da riconoscere dentro di s\'{e}: nessuno voleva essere come Achille, se ne riconosceva semmai il proprio potenziale di rabbia distruttiva. L'identificazione diretta Disney crea una forma di infantilismo di massa: adulti che continuano a cercare nella realt\`{a} i pattern narrativi dell'infanzia, che si aspettano finali da favola dalle proprie relazioni, e che interpretano la normale complessit\`{a} emotiva come una deviazione dalla norma.

\section{Il merchandising della moralit\`{a}}

Non si pu\`{o} parlare di Disney senza affrontare il denaro. La riduzione dei personaggi a immagini eterne e facilmente vendibili non \`{e} un effetto collaterale della creativit\`{a} Disney: ne \`{e} il motore principale.\footnote{Giroux, H. A. (1999). \textit{The Mouse that Roared}. Rowman e Littlefield, Lanham.} Belle diventa quella che legge, Mulan quella coraggiosa, Cenerentola quella servizievole: ogni principessa ridotta a uno slogan di marketing, riproducibile su tazze e magliette.

C'\`{e} un pensatore francese che ha analizzato questo esatto meccanismo di trasformazione con una precisione quasi chirurgica, molto prima che Disney lo perfezionasse su scala globale. Nel 1957 Roland Barthes pubblic\`{o} \textit{Mitologie}, una raccolta di brevi analisi in cui mostrava come la cultura popolare del suo tempo, dai detersivi al catch, dalle automobili alla bistecca con patatine, trasformasse sistematicamente l'ideologia in natura.\footnote{Barthes, R. (1957). \textit{Mythologies}. \'{E}ditions du Seuil, Parigi.} Il meccanismo che Barthes chiamava mito non \`{e} la favola nel senso comune, ma un particolare modo di parlare: un oggetto o un'immagine, costruiti storicamente e carichi di interessi ben precisi, vengono presentati come se fossero semplicemente cos\`{i} da sempre, ovvi, naturali, al di fuori di ogni discussione. Il mito, scriveva Barthes, ha il compito di far scomparire la storia dentro la natura, di trasformare una scelta in un'evidenza. \`{E} esattamente ci\`{o} che accade quando una principessa Disney smette di essere un personaggio costruito da sceneggiatori, animatori e reparti marketing, e diventa semplicemente quella che legge, un'etichetta naturale e immediata come il colore dei suoi capelli. Il merchandising della moralit\`{a} descritto sopra \`{e}, in termini barthesiani, mitologia allo stato puro applicata all'industria dell'intrattenimento: la complessit\`{a} psicologica umana viene prima ridotta a tratto unico, poi quel tratto viene naturalizzato al punto che nessuno si chiede pi\`{u} perch\'{e} il coraggio femminile debba assomigliare esattamente a Mulan. Quando riduciamo l'identit\`{a} a un marchio, insegniamo implicitamente che anche le persone reali sono, in fondo, marchi da costruire e vendere, e che l'autenticit\`{a} stessa \`{e} solo un'altra forma di performance.

\section{L'antidoto Miyazaki}

Fortunatamente esiste un antidoto quasi perfetto a questi meccanismi, e viene dal Giappone. Hayao Miyazaki non si \`{e} limitato a creare film alternativi: ha costruito un intero paradigma narrativo che rappresenta l'opposto filosofico del modello Disney.\footnote{Hern\'{a}ndez P\'{e}rez, M. (2016). Animation, branding and authorship in the construction of the anti-Disney ethos. \textit{Animation}, 11(1).} Quando visit\`{o} gli studi Disney nel 1999 per vedere \textit{Fantasia 2000}, la sua reazione fu brutalmente onesta: terribile, davvero terribile. Non era semplice rivalit\`{a} artistica: era il rifiuto di un'intera visione del mondo cinematografico.

Dove Disney utilizza la struttura occidentale a tre atti fondata su conflitto, climax e risoluzione, Miyazaki impiega il \textit{kish\=otenketsu} giapponese, una forma narrativa a quattro atti che non dipende necessariamente dal conflitto come motore. Questo approccio lascia spazio a ci\`{o} che chiama \textit{ma}, i momenti di vuoto contemplativo che, se riempiti di pensiero ed emozione, non fanno perdere il pubblico ma lo avvicinano. Miyazaki demolisce sistematicamente ogni meccanismo appena descritto. Contro la semplificazione morale, i suoi personaggi sono moralmente complessi: Lady Eboshi in \textit{Principessa Mononoke} distrugge la foresta sacra ma salva prostitute e malati, non \`{e} buona n\'{e} cattiva, \`{e} semplicemente umana. Contro la repressione emotiva, i suoi protagonisti si arrabbiano, dubitano, sbagliano: Chihiro ne \textit{La citt\`{a} incantata} \`{e} capricciosa e impaurita prima di crescere, e la sua crescita \`{e} faticosa quanto realistica, con quella catarsi che Disney tende a evitare. Contro il controllo narrativo, i suoi film sono pieni di simbolismo ambiguo, momenti contemplativi, significati che emergono lentamente, rispettando profondamente l'intelligenza del pubblico, anche di quello pi\`{u} piccolo. Contro l'identificazione diretta, i suoi personaggi non sono modelli da imitare ma specchi in cui riconoscere parti di s\'{e}: No Face non \`{e} un villain da sconfiggere, ma una rappresentazione della nostra stessa vuotezza consumistica, l'Ombra junghiana travestita da spirito senza volto. \`{E} l'anti Disney perfetto: profondo dove Disney \`{e} superficiale, complesso dove Disney semplifica, rispettoso dove Disney condiscende.

Miyazaki, per\`{o}, resta per sua stessa natura un'esperienza in parte elitaria: i suoi simbolismi richiedono un minimo di alfabetizzazione culturale, le sue pause contemplative possono apparire lente a bambini abituati a stimoli costanti. Un'alternativa pi\`{u} accessibile, ma altrettanto rivoluzionaria nella sostanza, \`{e} arrivata pi\`{u} di recente dall'Australia, sotto forma di una serie animata su una famiglia di cani. Quella serie ha saputo mantenere la semplicit\`{a} di linguaggio necessaria ai bambini pi\`{u} piccoli incorporando per\`{o} la stessa sofisticazione psicologica: genitori imperfetti che perdono la pazienza e poi chiedono scusa, emozioni difficili trattate come parte normale delle relazioni e non come segno di villania, temi delicati affrontati con sottigliezza invece che con moralismo esplicito. \`{E} una sorta di Miyazaki democratico: stessa filosofia narrativa, packaging pi\`{u} familiare, e la dimostrazione che accessibilit\`{a} e profondit\`{a} non sono affatto in contraddizione tra loro.

\section{La resistenza \`{e} possibile}

Cosa fare, allora? Abolire Disney dalle nostre vite, crescere i bambini senza le sue storie? Probabilmente no, e forse non sarebbe nemmeno desiderabile. Il problema non \`{e} Disney in s\'{e}, ma l'assenza di alternative, la mancanza di diversit\`{a} narrativa, l'egemonia culturale di un singolo paradigma di racconto.\footnote{Zipes, J. (1995). \textit{Breaking the Magic Spell}. University Press of Kentucky, Lexington.} La soluzione non \`{e} la censura, ma quella che potremmo chiamare alfabetizzazione narrativa: insegnare a riconoscere i pattern, a notare cosa manca in una storia troppo levigata, a cercare attivamente narrazioni pi\`{u} complesse e sfumate. \`{E} come insegnare a nuotare in un oceano di propaganda: non si pu\`{o} prosciugare l'oceano, ma si pu\`{o} imparare a non affogare.

Il successo internazionale sia di Miyazaki, premiato con un Oscar, sia della serie australiana appena descritta, diventata tra le pi\`{u} viste al mondo nella sua fascia, dimostra che esiste una fame diffusa di narrativa pi\`{u} sofisticata anche nell'intrattenimento familiare. Bettelheim ci ha insegnato perch\'{e} le storie difficili servano, Jung ci ha mostrato di quali materiali profondi siano fatte, Barthes ci ha svelato come l'industria le trasformi in merce naturalizzata, Aristotele ci ricorda a cosa serva davvero attraversare un'emozione invece di evitarla. Le storie che raccontiamo ai bambini, e che continuiamo silenziosamente a raccontarci da adulti, non sono mai innocenti: sono l'ultimo, pi\`{u} dolce veicolo attraverso cui i nostri bias pi\`{u} antichi continuano a viaggiare, generazione dopo generazione, travestiti da lieto fine.